%%Tema para beamer "Cuerna", versión 1.0
%%Desarrollado por Geri Morales (gerino.morales@im.unam.mx)
%%Imate Cuernavaca, UNAM. Agosto 2014.

\documentclass{beamer}
\usepackage[utf8]{inputenc}
\usepackage[T1]{fontenc} 
\usepackage[spanish]{babel} 
\usepackage{graphicx}
%%Se define el "environment" teorema
\newtheorem{teorema}{Teorema} 


%%Tema de beamer "Cuerna"
\usetheme{Cuerna} 
%\usetheme{Cuerna} 
%%Si se omite "[cuernavaca]" en éste comando, el logotipo se imprime sin la 
%%leyenda "Unidad Cuernavaca" en la parte inferior.

\usecolortheme{bluesimplex}
%Available color themes: default, bluesimplex, brick, lettuce

\logo{\includegraphics[width=1.7cm]{bluesimplex.png}}

\title{Título de la presentación}
\author{Nombre Apellido \\ %Nombre del autor 
	      \texttt{fulano@im.unam.mx}} %Dirección de correo electrónico (opcional, cuidado con la llave de cierre)
	      
\date{Congreso importante 2014} %Fecha o evento en que se presentará la plática

%\institute{Instituto de Matemáticas, unidad Cuernavaca} 
%%Instituto del ponenete, dado que el texto "Intituto de matemáticas" aparece
%%en el logo, parece redundante incluirlo además con éste comando.

\begin{document}

\begin{frame}

  \titlepage %Necesario para generar la portada
  
\end{frame}

%%La siguiente diapositiva es opcional, si se quiere la tabla de contenidos
%%Se sebe compilar dos veces el documento para que funcione
%--------------------------------------------------------------------------
\begin{frame}
\tableofcontents %Imprime la tabla de contenido
\end{frame}
%--------------------------------------------------------------------------

\section{Introducción} %%Título de la sección (Opcional)
\begin{frame}
  \frametitle{Título de la diapositiva}
  \framesubtitle{Subtítulo} %%Subtítulo de la diapositiva (opcional)
  Algo de texto en la parte superior
  
  \[
   \int_{0}^{\infty} \frac{5x^2}{\sqrt{a+b}}\, dx
   \]

  \begin{itemize}
    \item[\checkmark] Un elemento en la lista %%[\checkmark] muestra una palomita al inicio de la línea
    \item Segundo elemento
    \item Otro elemento
    \item Y finalmente, otro más
  \end{itemize}
\end{frame}

\begin{frame}{Otra diapositiva} %%Otra forma (más corta) de poner el título a la diapositiva
Texto adicional. No tiene un propósito en particular más que ocupar algo de espacio.
Texto adicional. No tiene un propósito en particular más que ocupar algo de espacio.
Texto adicional. No tiene un propósito en particular más que ocupar algo de espacio.
Texto adicional. No tiene un propósito en particular más que ocupar algo de espacio.
\( \sum_0^{\infty} a_i \) 
Texto adicional. No tiene un propósito en particular más que ocupar algo de espacio.
\end{frame}

\section{Segunda sección} %%Otra sección
\begin{frame}
    \frametitle{Un teorema}
    
    \begin{teorema} %%Uso del "environment" definido al inicio del documento.
        Los números primos son infinitos
    \end{teorema}

    Y algo de texto
\end{frame}

\end{document}